\documentclass[11pt,a4paper]{article}

\usepackage[utf8]{inputenc}
\usepackage{amsmath,amssymb}
\usepackage{booktabs,tabularx,longtable}
\usepackage{enumitem}
\usepackage{geometry}
\usepackage{hyperref}
\usepackage{xcolor}
\usepackage{listings}
\usepackage{fancyhdr}
\usepackage{titlesec}
\usepackage{tocloft}

\geometry{margin=2.5cm,top=3cm,bottom=3cm}

\hypersetup{
  colorlinks=true,
  linkcolor=blue!60!black,
  citecolor=blue!60!black,
  urlcolor=blue!60!black,
  pdftitle={ISLS Full-Stack Autonomy},
  pdfauthor={Sebastian Klemm},
}

\lstset{
  basicstyle=\ttfamily\small,
  breaklines=true,
  frame=single,
  backgroundcolor=\color{gray!8},
  keywordstyle=\color{blue!70!black}\bfseries,
  stringstyle=\color{red!60!black},
  commentstyle=\color{green!50!black}\itshape,
  numbers=left,
  numberstyle=\tiny\color{gray},
  numbersep=8pt,
  tabsize=4,
  showstringspaces=false,
}

\titleformat{\section}{\Large\bfseries}{\thesection}{1em}{}
\titleformat{\subsection}{\large\bfseries}{\thesubsection}{1em}{}
\titleformat{\subsubsection}{\normalsize\bfseries}{\thesubsubsection}{1em}{}

\pagestyle{fancy}
\fancyhf{}
\fancyhead[L]{\small ISLS Full-Stack Autonomy v1.0}
\fancyhead[R]{\small Klemm, 2026}
\fancyfoot[C]{\thepage}

\begin{document}

\begin{titlepage}
\begin{center}
\vspace*{2cm}

{\Huge\bfseries ISLS Full-Stack Autonomy}\\[0.8cm]
{\LARGE\bfseries From Application Description\\to Deployed System}\\[1.5cm]
{\Large Specification v1.0}\\[2cm]

{\large Sebastian Klemm}\\[0.3cm]
{\normalsize Independent Researcher, Germany}\\[0.5cm]
{\small March 2026}\\[2cm]

\begin{tabular}{ll}
\textbf{Repository:} & \texttt{graphity} \\
\textbf{Foundation:} & PSE (Post-Symbolic Engine) \\
\textbf{Eyes:} & Barbara (Code Topology, fused into ISLS) \\
\textbf{Hands:} & Die Schmiede (Language Generation, fused into ISLS) \\
\textbf{License:} & MIT \\
\end{tabular}

\vfill

{\small
This specification describes the fusion of PSE, Barbara, and Die Schmiede\\
into ISLS as a single recursive system that generates, analyzes, validates,\\
and deploys complete full-stack applications from a natural language description.\\[0.5cm]
All subsystems live in one repository: \texttt{graphity}.\\
No separate repos. No JSON bridges. No external dependencies.\\
One workspace. One binary. One \texttt{cargo run}.
}

\end{center}
\end{titlepage}

\tableofcontents
\newpage

% ============================================================
\section{Vision}

\subsection{The 3D Printer for Software}

A 3D printer takes a CAD model and produces a physical object.
Layer by layer, autonomously, without human intervention after
the print job starts.

ISLS takes an application description and produces a deployable
software system. Layer by layer---data models, database, services,
API, frontend, tests, deployment---autonomously, without human
intervention after the job starts.

\begin{lstlisting}[title=The Input]
isls forge --app "Warehouse Management System" \
           --constraints warehouse.toml \
           --output ./warehouse-system/
\end{lstlisting}

\begin{lstlisting}[title=The Output (10 minutes later)]
warehouse-system/
|-- backend/
|   |-- Cargo.toml
|   |-- src/
|   |   |-- main.rs              (Actix-web server, all routes wired)
|   |   |-- config.rs            (Environment config, DB connection)
|   |   |-- models/
|   |   |   |-- product.rs       (struct, DB mapping, validation)
|   |   |   |-- order.rs
|   |   |   |-- warehouse.rs
|   |   |   |-- user.rs
|   |   |-- services/
|   |   |   |-- inventory.rs     (business logic: stock, reorder)
|   |   |   |-- orders.rs        (create, fulfill, cancel)
|   |   |   |-- reporting.rs     (aggregation, export)
|   |   |   |-- auth.rs          (JWT, roles)
|   |   |-- api/
|   |   |   |-- products.rs      (CRUD endpoints)
|   |   |   |-- orders.rs
|   |   |   |-- reports.rs
|   |   |   |-- auth.rs
|   |   |-- database/
|   |   |   |-- migrations/      (SQL, versioned)
|   |   |   |-- queries.rs       (prepared statements)
|   |   |   |-- pool.rs          (connection pool)
|   |-- tests/
|   |   |-- api_tests.rs         (integration: all endpoints)
|   |   |-- service_tests.rs
|-- frontend/
|   |-- index.html
|   |-- src/
|   |   |-- app.js               (SPA, routing)
|   |   |-- pages/
|   |   |   |-- dashboard.js     (KPIs, charts)
|   |   |   |-- inventory.js     (table, search, CRUD)
|   |   |   |-- orders.js        (list, detail, status)
|   |   |   |-- reports.js       (date range, export)
|   |   |-- components/
|   |   |   |-- table.js         (sortable, filterable)
|   |   |   |-- form.js          (validation, submit)
|   |   |   |-- chart.js         (bar, line, pie)
|   |   |   |-- navbar.js
|   |   |-- api/
|   |   |   |-- client.js        (fetch wrapper, auth headers)
|   |-- style.css
|-- docker-compose.yml            (backend + postgres + frontend)
|-- README.md                     (setup, usage, API docs)
|-- evidence/
|   |-- generation_manifest.json  (what was generated, in what order)
|   |-- crystal_registry.json     (all crystals produced)
|   |-- constraint_proof.json     (SHA-256 chain proving constraints met)
|   |-- topology_report.json      (Barbara's analysis of generated code)
\end{lstlisting}

Everything compiles. Tests pass. Docker starts. The application works.
Generated autonomously. Verified topologically. Evidence-chained
cryptographically.

\subsection{Why This Requires Fusion}

Separate systems connected by JSON files cannot achieve this because:

\begin{enumerate}
\item The \textbf{planning layer} needs to understand code topology
  (Barbara) to decide what to generate next and in what order.
\item The \textbf{generation layer} needs structural blueprints
  (Schmiede concepts) to constrain the LLM oracle.
\item The \textbf{validation layer} needs to read the generated code
  back (Barbara) and verify it matches the plan.
\item The \textbf{integration layer} needs to compare modules
  topologically (Barbara) to verify interfaces match.
\item The \textbf{learning layer} needs to crystallize successful
  patterns (PSE) and use them to guide future generation.
\item All of this happens \textbf{within a single tick loop} ---
  latency of JSON file I/O between separate processes would kill
  the feedback speed.
\end{enumerate}

Therefore: PSE, Barbara, and Schmiede are \textbf{absorbed} into
ISLS. Not as dependencies. As integrated subsystems within the
same Rust workspace, sharing types, sharing memory, sharing the
tick loop.

% ============================================================
\section{Architecture}

\subsection{Subsystem Fusion Map}

\begin{center}
\small
\begin{tabular}{lll}
\toprule
\textbf{Source} & \textbf{Absorbed Into} & \textbf{What It Becomes} \\
\midrule
PSE pse-core & isls-engine & Engine orchestrator with full PSE pipeline \\
PSE pse-topology & isls-topology & Laplacian, Fiedler, Betti (already exists) \\
PSE pse-scale & isls-scale & Kuramoto synchronization (already exists) \\
PSE pse-memory & isls-memory & Pattern memory with cross-session persistence \\
PSE pse-navigator & isls-navigator & TRITON active exploration (already exists) \\
PSE pse-swarm & isls-swarm & Distributed consensus (already exists) \\
Barbara codex-parse & isls-reader & tree-sitter multi-language parser \\
Barbara codex-topo & isls-code-topo & Code-specific topology computation \\
Barbara codex-learn & isls-learner & Cross-language pattern accumulation \\
Schmiede schmiede-grammar & isls-langgen & Grammar generation from constraints \\
Schmiede schmiede-*gen & isls-langgen & Lexer/parser/typechecker/codegen generation \\
New & isls-planner & Architecture planning from description \\
New & isls-orchestrator & Multi-module generation sequencing \\
New & isls-integrator & Cross-module interface verification \\
New & isls-deployer & Docker/config/deployment generation \\
\bottomrule
\end{tabular}
\end{center}

Note: existing ISLS crates (isls-forge, isls-foundry, isls-oracle,
isls-cascade, isls-types, isls-graph, isls-evidence, isls-registry,
etc.) remain unchanged. The new crates extend them.

\subsection{The Full-Stack Pipeline}

\begin{lstlisting}[title=11-Stage Full-Stack Pipeline]
Stage 0:  DESCRIBE    Natural language -> structured AppSpec
Stage 1:  PLAN        AppSpec -> Architecture (modules, dependencies, order)
Stage 2:  MODEL       Architecture -> Data Models (structs, DB schema)
Stage 3:  PERSIST     Data Models -> Database Layer (migrations, queries)
Stage 4:  LOGIC       Architecture -> Service Layer (business logic)
Stage 5:  API         Services -> API Layer (endpoints, routes, auth)
Stage 6:  FRONTEND    API spec -> Frontend (pages, components, API client)
Stage 7:  TEST        All layers -> Integration Tests + E2E Tests
Stage 8:  DEPLOY      All layers -> Docker, configs, README
Stage 9:  VERIFY      Full system -> compile, test, topology check
Stage 10: LEARN       Results -> crystallize patterns, update memory
\end{lstlisting}

Each stage:
\begin{enumerate}
\item Checks the Blueprint Registry for known patterns (PSE Memory)
\item If found: generates from blueprint (fast, no Oracle needed)
\item If not found: constructs a constrained prompt, calls Oracle
\item Reads back the generated code (Barbara/isls-reader)
\item Verifies topology matches plan (isls-code-topo)
\item If mismatch: regenerates (up to 3 attempts)
\item If match: crystallizes the pattern, adds to registry
\item Passes output to next stage
\end{enumerate}

\subsection{Crate Map}

\begin{lstlisting}[language=bash,title=graphity/isls/ Workspace After Fusion]
isls/
|-- Cargo.toml
|-- crates/
|   |-- # === EXISTING (unchanged) ===
|   |-- isls-types/          (Crystal, Observation, Evidence)
|   |-- isls-graph/          (Observation graph)
|   |-- isls-extract/        (Pattern extraction)
|   |-- isls-cascade/        (8-gate adversarial validation)
|   |-- isls-evidence/       (SHA-256 evidence chains)
|   |-- isls-registry/       (Crystal registry)
|   |-- isls-topology/       (Laplacian, spectral decomposition)
|   |-- isls-scale/          (Kuramoto, multi-scale)
|   |-- isls-navigator/      (TRITON golden-angle spiral)
|   |-- isls-swarm/          (Multi-agent consensus)
|   |-- isls-constraint/     (DoF analysis)
|   |-- isls-capsule/        (AES-256-GCM encryption)
|   |-- isls-store/          (SQLite persistence)
|   |-- isls-scheduler/      (Tick-based scheduling)
|   |-- isls-replay/         (Deterministic replay)
|   |-- isls-manifest/       (Manifest construction)
|   |-- isls-pmhd/           (Adversarial drill)
|   |-- isls-forge/          (Code skeleton generation)
|   |-- isls-foundry/        (Compile-test-fix loop)
|   |-- isls-oracle/         (LLM prompt construction)
|   |-- isls-templates/      (Architecture templates)
|   |-- isls-compose/        (Code composition)
|   |-- isls-artifact-ir/    (Code intermediate representation)
|   |-- isls-babylon/        (Multi-language generation - legacy)
|   |-- isls-agent/          (No-code operator)
|   |-- isls-studio/         (Web UI)
|   |-- isls-gateway/        (HTTP API)
|   |-- isls-cli/            (CLI)
|   |--
|   |-- # === NEW: Fused from Barbara ===
|   |-- isls-reader/         (tree-sitter parsing, 8+ languages)
|   |-- isls-code-topo/      (Code topology computation)
|   |-- isls-learner/        (Cross-language pattern accumulation)
|   |-- isls-memory/         (Persistent pattern memory)
|   |--
|   |-- # === NEW: Fused from Schmiede ===
|   |-- isls-langgen/        (Grammar + lexer + parser + typechecker + codegen generation)
|   |--
|   |-- # === NEW: Full-Stack Pipeline ===
|   |-- isls-planner/        (AppSpec -> Architecture)
|   |-- isls-orchestrator/   (Multi-stage generation sequencing)
|   |-- isls-integrator/     (Cross-module verification)
|   |-- isls-deployer/       (Docker, configs, README generation)
|   |-- isls-blueprint/      (Blueprint registry + matching)
\end{lstlisting}

% ============================================================
\section{Stage Specifications}

\subsection{Stage 0: DESCRIBE}

Input: natural language application description + constraint TOML.

\begin{lstlisting}[title=Constraint TOML (application-level)]
[app]
name = "warehouse-system"
description = "Warehouse management with inventory, orders, reporting"

[app.modules]
inventory = "Product catalog, stock levels, reorder alerts"
orders = "Order creation, fulfillment, cancellation, tracking"
reporting = "Daily/weekly/monthly reports, KPI dashboard"
auth = "JWT authentication, role-based access (admin, operator)"

[backend]
language = "rust"
framework = "actix-web"
database = "postgresql"
auth_method = "jwt"

[frontend]
type = "spa"            # single-page application
framework = "vanilla"   # no framework, vanilla JS
styling = "minimal"     # clean CSS, no Tailwind/Bootstrap

[deployment]
containerized = true
compose = true

[constraints]
max_crates = 1          # monolith backend (single crate)
test_coverage = "integration"  # at minimum integration tests
evidence_chain = true   # full SHA-256 evidence on every artifact
\end{lstlisting}

Output: \texttt{AppSpec} --- a structured, validated application specification.

\begin{lstlisting}[language=C,title=AppSpec]
pub struct AppSpec {
    pub name: String,
    pub description: String,
    pub modules: Vec<ModuleSpec>,
    pub backend: BackendSpec,
    pub frontend: FrontendSpec,
    pub deployment: DeploymentSpec,
    pub constraints: AppConstraints,
}

pub struct ModuleSpec {
    pub name: String,
    pub description: String,
    pub entities: Vec<String>,     // inferred: "Product", "Order", etc.
    pub operations: Vec<String>,   // inferred: "create", "list", "update", etc.
    pub dependencies: Vec<String>, // inferred: orders depends on inventory
}

pub struct BackendSpec {
    pub language: String,
    pub framework: String,
    pub database: String,
    pub auth: AuthSpec,
}

pub struct FrontendSpec {
    pub app_type: String,
    pub framework: String,
    pub styling: String,
    pub pages: Vec<PageSpec>,      // inferred from modules
}

pub struct PageSpec {
    pub name: String,
    pub route: String,
    pub components: Vec<String>,   // inferred: "table", "form", "chart"
}
\end{lstlisting}

\textbf{Implementation:} Parse the TOML. Infer entities from module
descriptions (nouns become entities). Infer operations (verbs become
CRUD operations). Infer dependencies (if module A references entities
from module B, A depends on B). Infer frontend pages (one page per
module, plus dashboard).

No LLM needed for Stage 0 --- this is deterministic parsing and
inference.

\subsection{Stage 1: PLAN}

Input: \texttt{AppSpec}.
Output: \texttt{Architecture} --- the generation plan.

\begin{lstlisting}[language=C,title=Architecture]
pub struct Architecture {
    pub layers: Vec<Layer>,
    pub generation_order: Vec<GenerationStep>,
    pub interfaces: Vec<Interface>,
    pub estimated_files: usize,
    pub estimated_loc: usize,
}

pub struct Layer {
    pub name: String,         // "models", "database", "services", "api", "frontend"
    pub components: Vec<Component>,
}

pub struct Component {
    pub name: String,         // "product_model", "inventory_service", "products_api"
    pub layer: String,
    pub file_path: String,    // "src/models/product.rs"
    pub depends_on: Vec<String>,
    pub blueprint: Option<BlueprintId>,  // matched from registry
    pub estimated_loc: usize,
}

pub struct GenerationStep {
    pub order: usize,         // 1, 2, 3, ...
    pub component: String,
    pub reason: String,       // "models first because services depend on them"
    pub can_parallel: Vec<String>,  // components that can be generated in parallel
}

pub struct Interface {
    pub from: String,         // "inventory_service"
    pub to: String,           // "products_api"
    pub interface_type: InterfaceType,
    pub contract: String,     // "fn list_products() -> Vec<Product>"
}

pub enum InterfaceType {
    FunctionCall,
    HttpEndpoint,
    DatabaseQuery,
    EventEmission,
}
\end{lstlisting}

\textbf{Generation order} follows dependency topology:
\begin{enumerate}
\item Data models (no dependencies)
\item Database migrations + queries (depends on models)
\item Services (depends on models + database)
\item API endpoints (depends on services)
\item Frontend pages (depends on API)
\item Tests (depends on all)
\item Deployment (depends on all)
\end{enumerate}

The Planner uses the Blueprint Registry to match components against
known patterns. If ``inventory\_service'' matches a known
``CRUDService'' blueprint, the planner records this --- the Forge
will use the blueprint for guided generation.

\textbf{Implementation:} Deterministic. No LLM. The Planner takes
the AppSpec's modules, creates one component per entity per layer,
wires dependencies, and determines generation order via topological
sort.

\subsection{Stage 2: MODEL}

Input: Architecture (model layer components).
Output: Rust structs + SQL schema.

For each entity in each module, generate:
\begin{itemize}
\item A Rust struct with serde derives
\item Validation methods
\item SQL CREATE TABLE statement
\item FromRow / ToRow conversion
\end{itemize}

\textbf{Example output:}
\begin{lstlisting}[language=C]
// src/models/product.rs
#[derive(Debug, Clone, Serialize, Deserialize)]
pub struct Product {
    pub id: i64,
    pub name: String,
    pub sku: String,
    pub quantity: i32,
    pub reorder_level: i32,
    pub price_cents: i64,
    pub warehouse_id: i64,
    pub created_at: DateTime<Utc>,
    pub updated_at: DateTime<Utc>,
}

impl Product {
    pub fn validate(&self) -> Result<(), ValidationError> {
        if self.name.is_empty() { return Err(ValidationError::EmptyField("name")); }
        if self.quantity < 0 { return Err(ValidationError::NegativeQuantity); }
        if self.price_cents < 0 { return Err(ValidationError::NegativePrice); }
        Ok(())
    }
}
\end{lstlisting}

\textbf{Implementation:} Template-based for simple CRUD entities (no
Oracle needed). LLM Oracle only for complex domain logic that can't
be inferred from the description.

\subsection{Stage 3: PERSIST}

Input: Data models.
Output: SQL migrations + Rust database access layer.

Generate:
\begin{itemize}
\item \texttt{migrations/001\_initial.sql} --- CREATE TABLE for each model
\item \texttt{database/queries.rs} --- CRUD queries (sqlx or raw SQL)
\item \texttt{database/pool.rs} --- connection pool setup
\end{itemize}

Template-driven. No Oracle for standard CRUD.

\subsection{Stage 4: LOGIC}

Input: Architecture (service layer components) + models + database.
Output: Service modules with business logic.

This is where the Oracle is needed. Business logic can't be fully
templated: ``reorder when stock below threshold and no pending
orders exist'' requires understanding of domain semantics.

The Forge:
\begin{enumerate}
\item Checks Blueprint Registry for matching service pattern
\item Constructs prompt with: model structs (context), database
  queries (available), service description (from AppSpec),
  blueprint constraints (if matched)
\item Calls Oracle
\item Reads back generated code with isls-reader (Barbara)
\item Verifies topology matches blueprint
\item If mismatch: regenerates
\item If match: crystallizes, moves to next component
\end{enumerate}

\subsection{Stage 5: API}

Input: Services.
Output: HTTP endpoint handlers.

Mostly template-driven for CRUD:
\begin{lstlisting}[language=C]
// Template: for each entity with CRUD operations
#[get("/api/{entity_plural}")]
async fn list_{entity}s(db: Data<Pool>) -> impl Responder {
    let items = db.query("{entity}_list_all").await?;
    HttpResponse::Ok().json(items)
}

#[post("/api/{entity_plural}")]
async fn create_{entity}(db: Data<Pool>, body: Json<Create{Entity}>) -> impl Responder {
    body.validate()?;
    let item = db.query("{entity}_insert", &body).await?;
    HttpResponse::Created().json(item)
}
// ... update, delete, get_by_id
\end{lstlisting}

Non-CRUD endpoints (reporting aggregations, complex queries) go
through the Oracle with blueprint guidance.

\subsection{Stage 6: FRONTEND}

Input: API specification (all endpoints, request/response types).
Output: Single-page application (HTML + JS + CSS).

For each module/page:
\begin{itemize}
\item HTML page with table (list view) and form (create/edit)
\item JavaScript: fetch from API, render, handle events
\item Navigation between pages
\item Dashboard with KPI cards and charts
\end{itemize}

Template-driven for standard CRUD UI. Charts and complex
interactions may need Oracle.

\subsection{Stage 7: TEST}

Input: All generated layers.
Output: Integration tests + E2E tests.

\begin{itemize}
\item For each API endpoint: test with valid input (expect 200),
  test with invalid input (expect 400/422), test without auth
  (expect 401)
\item For each service: test core business logic
\item E2E: start server, create entity, read back, update, delete,
  verify
\end{itemize}

\textbf{Validation:} isls-foundry compiles and runs the tests.
If tests fail, the Foundry enters its compile-test-fix loop:
\begin{enumerate}
\item Read test failure message
\item Identify which component is broken
\item Regenerate that component (with failure context in prompt)
\item Recompile, retest
\item Up to 3 fix attempts per component
\end{enumerate}

\subsection{Stage 8: DEPLOY}

Input: All layers.
Output: docker-compose.yml, Dockerfiles, .env.example, README.md.

Template-driven. No Oracle needed.

\begin{lstlisting}[language=bash,title=docker-compose.yml (generated)]
version: '3.8'
services:
  db:
    image: postgres:16
    environment:
      POSTGRES_DB: warehouse
      POSTGRES_PASSWORD: ${DB_PASSWORD}
    volumes:
      - pgdata:/var/lib/postgresql/data
      - ./database/migrations:/docker-entrypoint-initdb.d
  backend:
    build: ./backend
    ports: ["8080:8080"]
    environment:
      DATABASE_URL: postgres://postgres:${DB_PASSWORD}@db/warehouse
    depends_on: [db]
  frontend:
    build: ./frontend
    ports: ["3000:80"]
volumes:
  pgdata:
\end{lstlisting}

\subsection{Stage 9: VERIFY}

Input: Complete generated system.
Output: VerificationReport.

\begin{enumerate}
\item \texttt{cargo build} the backend --- must compile
\item \texttt{cargo test} --- all tests must pass
\item isls-reader (Barbara) reads ALL generated Rust code
\item isls-code-topo computes topology of each module
\item isls-integrator verifies cross-module interfaces:
  \begin{itemize}
  \item Every API endpoint has a corresponding service function
  \item Every service function uses only existing database queries
  \item Every frontend API call hits an existing endpoint
  \item Every model used in a service is defined in models/
  \end{itemize}
\item Evidence chain verification: every artifact has SHA-256 hash,
  chain is unbroken from Stage 0 to Stage 9
\end{enumerate}

\subsection{Stage 10: LEARN}

Input: Verification results.
Output: Updated Blueprint Registry.

For every component that was:
\begin{itemize}
\item Generated on first attempt (no fix needed) $\rightarrow$ high-confidence
  blueprint (0.9)
\item Generated after 1 fix $\rightarrow$ medium-confidence blueprint (0.7)
\item Generated after 2--3 fixes $\rightarrow$ low-confidence blueprint (0.5)
\item Failed all attempts $\rightarrow$ negative pattern (``don't do this'')
\end{itemize}

The blueprints are stored in the ISLS crystal registry. Next time
ISLS generates a similar application, these blueprints guide the
Forge. Over time, common patterns (CRUD service, REST endpoint,
login flow) become fully templated --- no Oracle needed.

This is Mining: the system runs, generates, learns, and the next
run is faster and better. Indefinitely.

% ============================================================
\section{New Crate Specifications}

\subsection{isls-reader (from Barbara codex-parse)}

\begin{lstlisting}[language=C]
/// Parse any source file via tree-sitter.
/// Supports: Rust, Python, JavaScript, TypeScript, Go, C, Java, SQL, HTML, CSS.
pub fn parse_file(path: &Path) -> Result<Vec<CodeObservation>>;

/// Parse a string with explicit language.
pub fn parse_string(source: &str, language: &str) -> Result<Vec<CodeObservation>>;

/// Parse all source files in a directory (recursive).
pub fn parse_directory(dir: &Path) -> Result<WorkspaceAnalysis>;
\end{lstlisting}

Port from Barbara's codex-parse. Include tree-sitter grammars for
all languages that ISLS might generate: Rust (primary), JavaScript,
HTML, CSS, SQL, TOML, Dockerfile.

\subsection{isls-code-topo (from Barbara codex-topo)}

\begin{lstlisting}[language=C]
/// Compute topological signature of code.
pub fn compute_code_topology(observations: &[CodeObservation]) -> CodeTopology;

/// Compare two code topologies.
pub fn topology_similarity(a: &CodeTopology, b: &CodeTopology) -> f64;
\end{lstlisting}

Port from Barbara's codex-topo with the hardened weights (95\%
recognition, 0.88 threshold).

\subsection{isls-planner}

\begin{lstlisting}[language=C]
/// Create an architecture plan from an AppSpec.
pub fn plan(
    spec: &AppSpec,
    blueprints: &BlueprintRegistry,
) -> Result<Architecture>;

/// Determine generation order via topological sort.
pub fn generation_order(arch: &Architecture) -> Vec<GenerationStep>;

/// Estimate total LOC and generation time.
pub fn estimate(arch: &Architecture) -> Estimate;
\end{lstlisting}

\subsection{isls-orchestrator}

\begin{lstlisting}[language=C]
/// The main loop. Drives the entire 11-stage pipeline.
pub struct Orchestrator {
    planner: Planner,
    forge: Forge,
    foundry: Foundry,
    reader: Reader,
    topo: CodeTopoAnalyzer,
    integrator: Integrator,
    deployer: Deployer,
    blueprints: BlueprintRegistry,
    memory: PatternMemory,
}

impl Orchestrator {
    /// Run the full pipeline from AppSpec to deployed system.
    pub fn run(&mut self, spec: AppSpec, output_dir: &Path) -> Result<RunReport>;

    /// Run a single stage.
    pub fn run_stage(&mut self, stage: Stage) -> Result<StageResult>;

    /// Resume a partially completed run.
    pub fn resume(&mut self, output_dir: &Path) -> Result<RunReport>;
}

pub struct RunReport {
    pub stages_completed: Vec<StageResult>,
    pub total_files_generated: usize,
    pub total_loc: usize,
    pub oracle_calls: usize,
    pub blueprint_hits: usize,
    pub blueprint_hit_rate: f64,
    pub compile_success: bool,
    pub tests_passed: usize,
    pub tests_failed: usize,
    pub crystals_created: usize,
    pub total_time_secs: f64,
    pub evidence_chain_valid: bool,
}

pub struct StageResult {
    pub stage: Stage,
    pub components_generated: usize,
    pub oracle_calls: usize,
    pub retries: usize,
    pub time_secs: f64,
    pub success: bool,
}
\end{lstlisting}

\subsection{isls-integrator}

\begin{lstlisting}[language=C]
/// Verify cross-module interfaces.
pub fn verify_integration(
    output_dir: &Path,
    architecture: &Architecture,
) -> Result<IntegrationReport>;

pub struct IntegrationReport {
    pub interfaces_checked: usize,
    pub interfaces_valid: usize,
    pub mismatches: Vec<InterfaceMismatch>,
    pub all_valid: bool,
}

pub struct InterfaceMismatch {
    pub from: String,
    pub to: String,
    pub expected: String,
    pub found: String,
    pub severity: Severity,
}
\end{lstlisting}

The integrator uses isls-reader to parse the generated code and
checks:
\begin{itemize}
\item Every function called exists in the target module
\item Every struct used is defined in the imported module
\item API route strings match frontend fetch URLs
\item SQL column names match Rust struct field names
\end{itemize}

\subsection{isls-deployer}

\begin{lstlisting}[language=C]
/// Generate deployment artifacts.
pub fn generate_deployment(
    architecture: &Architecture,
    spec: &AppSpec,
    output_dir: &Path,
) -> Result<Vec<PathBuf>>;
\end{lstlisting}

Generates: Dockerfile(s), docker-compose.yml, .env.example,
.gitignore, README.md. Template-driven, no Oracle.

\subsection{isls-blueprint (simplified from earlier spec)}

\begin{lstlisting}[language=C]
/// Blueprint registry for structural patterns.
pub struct BlueprintRegistry { ... }

impl BlueprintRegistry {
    pub fn find_match(&self, request: &GenerationRequest) -> Option<&Blueprint>;
    pub fn add(&mut self, blueprint: Blueprint);
    pub fn record_usage(&mut self, id: &str, success: bool);
    pub fn update_confidences(&mut self);
    pub fn save(&self, path: &Path) -> Result<()>;
    pub fn load(path: &Path) -> Result<Self>;
}
\end{lstlisting}

% ============================================================
\section{Implementation Strategy}

\subsection{Porting vs. Rewriting}

Barbara's codex-parse, codex-topo, and codex-learn are \textbf{ported}
into ISLS (isls-reader, isls-code-topo, isls-learner), not rewritten.
Copy the code, rename the crate, adjust import paths, verify tests pass.

Schmiede's grammar/lexgen/parsegen are \textbf{condensed} into a single
crate isls-langgen. The language generation capability remains
available but is not in the critical path for Full-Stack generation
(which generates Rust/JS/SQL, not custom languages).

\subsection{What Requires the LLM Oracle}

\begin{center}
\small
\begin{tabular}{lll}
\toprule
\textbf{Stage} & \textbf{Oracle Needed?} & \textbf{Reason} \\
\midrule
0 DESCRIBE & No & TOML parsing, deterministic \\
1 PLAN & No & Topological sort, deterministic \\
2 MODEL & Rarely & Templates for CRUD, Oracle for complex domain \\
3 PERSIST & No & SQL templates \\
4 LOGIC & \textbf{Yes} & Business logic requires semantic understanding \\
5 API & Rarely & Templates for CRUD endpoints \\
6 FRONTEND & Sometimes & Templates for standard UI, Oracle for complex \\
7 TEST & Sometimes & Templates for CRUD tests, Oracle for edge cases \\
8 DEPLOY & No & Templates \\
9 VERIFY & No & Compilation + topology check \\
10 LEARN & No & Pattern extraction, deterministic \\
\bottomrule
\end{tabular}
\end{center}

Of 11 stages, only Stage 4 (business logic) consistently needs the
Oracle. Everything else is template-driven or deterministic.
This means: \textbf{most of the application is generated without LLM
calls.} The Oracle is a specialist consulted for domain logic, not
a generalist that writes everything.

Over time, as the Blueprint Registry accumulates patterns from
successful generations, even Stage 4 calls decrease. The system
converges toward fully offline generation.

\subsection{Execution Order for Implementation}

\begin{enumerate}
\item \textbf{Port Barbara into ISLS:} isls-reader, isls-code-topo,
  isls-learner. Run Barbara's tests in the ISLS workspace. All must
  pass. This is mechanical: copy, rename, adjust paths.

\item \textbf{isls-blueprint:} Blueprint types, registry, matching.
  Test in isolation.

\item \textbf{isls-planner:} AppSpec parsing, architecture generation,
  topological sort. Test with the warehouse example.

\item \textbf{isls-orchestrator:} The main loop. Wire planner $\to$
  forge $\to$ foundry $\to$ reader $\to$ integrator. Initially with
  mock generation (templates only, no Oracle).

\item \textbf{Template generation} for Stages 2, 3, 5, 6, 8.
  These produce real Rust/JS/SQL code from templates.

\item \textbf{Forge enhancement:} Blueprint matching before Oracle
  call. Guided prompt construction. Post-generation topology check.

\item \textbf{isls-integrator:} Cross-module interface verification.

\item \textbf{isls-deployer:} Docker and config generation.

\item \textbf{Stage 9 VERIFY:} Full system compilation and test.

\item \textbf{Stage 10 LEARN:} Crystallization of successful patterns.

\item \textbf{End-to-end test:} warehouse example from TOML to
  compiled, tested, deployable system.
\end{enumerate}

% ============================================================
\section{The Warehouse Example}

The specification includes one complete example: a Warehouse
Management System. This serves as the integration test.

\subsection{Input}

The \texttt{warehouse.toml} constraint file (shown in Section 3.1).

\subsection{Expected Output Metrics}

\begin{center}
\begin{tabular}{lr}
\toprule
\textbf{Metric} & \textbf{Expected} \\
\midrule
Total files generated & 25--40 \\
Total LOC & 2,000--4,000 \\
Backend Rust files & 12--18 \\
Frontend JS/HTML/CSS files & 8--12 \\
SQL migration files & 1--3 \\
Config/deploy files & 3--5 \\
Test files & 2--4 \\
Oracle calls & 5--15 (mostly Stage 4) \\
Blueprint hits & depends on registry state \\
\texttt{cargo build} & PASS \\
\texttt{cargo test} & $\geq 80\%$ pass rate \\
Evidence chain & valid (all SHA-256 verified) \\
Total generation time & $< 15$ minutes \\
\bottomrule
\end{tabular}
\end{center}

\subsection{Verification Commands}

\begin{lstlisting}[language=bash]
# Generate
isls forge --app "Warehouse Management System" \
           --constraints warehouse.toml \
           --output ./output/warehouse

# Verify backend compiles
cd output/warehouse/backend && cargo build && cargo test

# Verify frontend serves
cd ../frontend && python3 -m http.server 3000 &

# Verify Docker
cd .. && docker-compose up -d

# Check evidence
cat evidence/constraint_proof.json | python3 -m json.tool | head
\end{lstlisting}

% ============================================================
\section{Success Criteria}

\begin{enumerate}
\item \texttt{cargo build --workspace} compiles with zero errors
  (all existing 446+ tests plus new tests).
\item Barbara's ported tests pass in the ISLS workspace
  (95\% cross-language recognition preserved).
\item \texttt{isls forge --constraints warehouse.toml} produces
  a complete output directory.
\item The generated backend compiles with \texttt{cargo build}.
\item At least 5 generated test cases pass.
\item The generated \texttt{docker-compose.yml} is syntactically valid.
\item The evidence chain from Stage 0 to Stage 9 is unbroken
  (SHA-256 verification passes).
\item Blueprint Registry has $\geq 5$ entries after one full run.
\item A second run of the same spec is measurably faster (blueprint
  hits reduce Oracle calls).
\end{enumerate}

% ============================================================
\section{Implementation Constraints for Claude Code}

\begin{enumerate}
\item \textbf{Working repository: graphity.} All code goes into
  the \texttt{isls/} subdirectory workspace.

\item \textbf{Barbara and Schmiede code must be cloned/ported.}
  Clone the barbara repository, copy the relevant crate source
  into isls/ under new names (isls-reader, isls-code-topo,
  isls-learner, isls-langgen). Adjust all import paths. Run
  their tests.

\item \textbf{tree-sitter dependency:} isls-reader needs tree-sitter
  and language grammars. Add to workspace dependencies. Requires
  C compiler (\texttt{gcc}).

\item \textbf{Template engine:} Use \texttt{tera} for code generation
  templates (Rust, JS, SQL, HTML, CSS, Docker, TOML).

\item \textbf{All 446 existing ISLS tests must pass.} Zero regressions.

\item \textbf{Sequential implementation.} Port Barbara first. Then
  blueprint. Then planner. Then orchestrator. Then templates. Then
  verify end-to-end.

\item \textbf{The warehouse end-to-end test is the most important
  deliverable.} If time runs out, skip isls-langgen (Schmiede port),
  skip isls-deployer (Docker), but ensure the backend generates,
  compiles, and has at least some passing tests.

\item \textbf{Mock Oracle.} If no LLM API key is available, implement
  a mock Oracle that returns template-based code for known component
  types (CRUD service, REST endpoint, model). The mock must be good
  enough to produce compilable code.

\item \textbf{No \texttt{unwrap()} in library code.}

\item \textbf{Every public item documented.}

\item \textbf{MIT License. Copyright (c) 2026 Sebastian Klemm.}
\end{enumerate}

\vfill
\noindent\rule{\textwidth}{0.4pt}
\begin{center}
\small
ISLS Full-Stack Autonomy v1.0 --- Sebastian Klemm --- March 2026\\
PSE + Barbara + Die Schmiede $\rightarrow$ fused into ISLS\\
One repository. One workspace. One \texttt{cargo run}.\\
\url{https://github.com/LashSesh/graphity}
\end{center}

\end{document}
